\chapter{Nondimensionalization and characteristic parameters}
\label{ch:dimensionless}

The governing groups are obtained by substituting the scales of the problem into the field
equations of Chapter~\ref{ch:equations}: length $H_0$, velocity $U$, pressure $\mu U/H_0$ (with the
equivalent inertial statement noted where used), temperature differences $\Delta T^*=q^{''}H_0/k_f$,
stress $E\alpha\Delta T^*$, displacement $\alpha\Delta T^*H_0$. No group is introduced
independently of the equations.

\textbf{Length-basis and friction-factor convention.} Two Reynolds numbers appear in the channel
literature: $\mathrm{Re}_{H_0}=\rho UH_0/\mu$ and $\mathrm{Re}_{D_h}=\rho U(2H_0)/\mu=2\mathrm{Re}_{H_0}$,
with $D_h=2H_0$ for the slit. The Darcy friction factor is defined by
$\Delta p=f_D\,(L/D_h)\,\rho U^2/2$, and for the slit $f_D\,\mathrm{Re}_{D_h}=96$
\cite{shah1978laminar}; the numerator differs by a factor four from the Fanning convention
($f_F\,\mathrm{Re}_{D_h}=24$). Throughout this document the $D_h$ basis and the Darcy convention
are used, and the basis is required to be printed on every computed output (Appendix~B, register
D-6).

\section{Groups from the fluid equations}
Substituting the scales into \eqref{eq:momentum}: the inertia-to-viscous ratio
\begin{equation}
\mathrm{Re}_{D_h}=\frac{\rho U(2H_0)}{\mu}\quad(=50\text{ at baseline}),
\end{equation}
and from \eqref{eq:fluidenergy} the advection-to-conduction ratio
\begin{equation}
\mathrm{Pr}=\frac{\mu c_p}{k_f}=6.97,\qquad \mathrm{Pe}=\mathrm{Re}_{D_h}\mathrm{Pr}=348 .
\end{equation}
Were the Boussinesq term retained, the ratio of buoyant to inertial forcing would be measured by
the Grashof and Richardson numbers
\begin{equation}
\mathrm{Gr}=\frac{g\beta_f\Delta T^*H_0^3}{\nu_{\mathrm{kin}}^2},\qquad
\mathrm{Ri}=\frac{\mathrm{Gr}}{\mathrm{Re}_{H_0}^2}
=\frac{g\beta_f\Delta T^*H_0}{U^2},
\label{eq:grri}
\end{equation}
where the second form makes the length basis explicit: both numerator and denominator of Ri must
carry the \emph{same} length, so with the Grashof number built on $H_0$ and the temperature scale
$\Delta T^*=q^{''}H_0/k_f$ (itself fixed by $H_0$), the consistent Reynolds number is
$\mathrm{Re}_{H_0}=\rho UH_0/\mu$, not $\mathrm{Re}_{D_h}$; equivalently, on the $D_h$ basis both
groups change ($\mathrm{Gr}_{D_h}=8\,\mathrm{Gr}_{H_0}$) and Ri changes by the factor
$D_h/H_0=2$. The two conventions differ by a factor of two and both are quoted in
Chapter~\ref{ch:baseline}. Finally, the dissipation group
\begin{equation}
\mathrm{Br}=\frac{\mu U^2}{k_f\Delta T^*}
\end{equation}
controls Assumption~A7. All four are evaluated in full arithmetic detail in
Section~\ref{sec:baselinecalc}; an earlier draft divided $\mathrm{Gr}(H_0)$ by
$\mathrm{Re}_{D_h}^2$ and mis-transcribed the powers of ten, values now corrected
(Appendix~B, register D-14).

\section{Groups from the solid and interface equations}
Continuity of the interface flux with the common scale $\Delta T^*$ brings in the conductivity
ratio and the two geometric ratios
\begin{equation}
R_k=\frac{k_s}{k_f}=0.42,\qquad D_r=\frac{d}{H_0}=6,\qquad \Lambda=\frac{L}{H_0}=60 ,
\end{equation}
written $R_k$ and $D_r$ to avoid collision with the bulk modulus $K$ and the rate-of-deformation
tensor $\mathbf D$ (Appendix~B, register D-18). The through-wall temperature drop expressed in
units of $\Delta T^*$ is
$\Delta T_t/\Delta T^*=D_r/R_k=14.4$: essentially the entire thermal gradient of the coupled
problem lives inside the wall, which is the quantitative basis of Assumption~A2's
self-consistency check and the reason the solid sees a much larger $\theta$ than the fluid ever
will.

The characteristic result of the model is that the relative thermal bowing of the channel,
from \eqref{eq:kappa}, \eqref{eq:wmax}, and $\Delta T_t=q^{''}d/k_s$, collapses to
\begin{equation}
\frac{w_{\max}^{\mathrm{lin}}}{H_0}=\frac{(1+\nu)\,\Sigma_\kappa\,\Lambda^2}{8},
\qquad\text{with}\qquad
\Sigma_\kappa=\alpha\,\frac{q^{''}H_0}{k_s}
\label{eq:relbow}
\end{equation}
a dimensionless curvature-strain number. At baseline $(1+\nu)\Sigma_\kappa=7\times10^{-4}$ and
$w_{\max}^{\mathrm{lin}}/H_0=0.315$; at the baseline membrane force the frozen-force
beam--column form of Section~\ref{sec:deriv:beamcolumn} gives the adopted
$w_{\max}/H_0=0.438$ (self-consistent comparison $0.428$), and \eqref{eq:relbow} is read as
the linear limit of that relation. Equation \eqref{eq:relbow} assembles geometry, material, and
load into a single ordinate and is the ordering parameter of the regime analysis in
Chapter~\ref{ch:baseline} and of any future parameter study.

The ratios governing the mechanical regime are the deformation ratio of
Section~\ref{sec:deriv:pressure},
$\Gamma_p=w_{\max}/w_p$, \eqref{eq:gammap}, and the buckling sub-criticality margin
\begin{equation}
M_b=\frac{\bar T_{\mathrm{cr}}}{\bar\theta}\qquad(\text{required } M_b\ge2) ,
\end{equation}
with $\bar T_{\mathrm{cr}}$ of \eqref{eq:tcr} and $\bar\theta$ the mean wall temperature rise.

\section{Admissibility bounds on the regime}
Any reported operating point of the model must lie inside the following bounds, all verified
arithmetic at the baseline in Chapter~\ref{ch:baseline}:
\begin{itemize}\itemsep1pt
\item[A-1] $\mathrm{Re}_{D_h}\le100$ (laminar-design headroom, A6);
\item[A-2] $M_b\ge2$ (sub-critical pinned supports, Section~\ref{sec:deriv:buckling});
\item[A-3] kinematics bounds (von K\'arm\'an validity, A4 and
Section~\ref{sec:deriv:beamcolumn}): bending fibre strain
$(\kappa_{\mathrm{geo}}-\kappa_{\mathrm{th}})d/2\le5\times10^{-3}$, total mechanical surface
strain $\le10^{-2}$, arc-stretching strain $\le10^{-3}$, maximum slope $\le0.1$~rad, plus the
mesh-generation bound $|w|/H_0\le0.5$ with positive element Jacobians;
\item[A-4] fluid-property drift below $\sim10\%$ over the computed fluid-side temperature range (A2);
\item[A-5] $\Gamma_p\ge5$ for a case to be classified as thermally driven, evaluated on the
\emph{like-for-like} definition of \eqref{eq:gammap} (the same compressive force amplifying both
responses, $1.73\times10^3$ at the baseline); the mixed value of \eqref{eq:gammap} is reported
alongside and is larger, so the criterion is met under either definition (register D-29);
\item[A-6] computed wall stress $\sigma_{\max}\le0.2\,\sigma_y$, with $\sigma_y$ from the frozen
datasheet band (design cap with margin $\ge5$; register D-19 restates the earlier ambiguous
wording).
\end{itemize}

\section{Dependence on yet-unfrozen data}
Only temperature differences appear in the equations, so $T_{\mathrm{ref}}$ and $T_{\mathrm{in}}$
enter the solution exclusively through property values; both are to be fixed at the property-data
freeze (Appendix~B, C-6), as is the value of $\beta_f$ entering $\mathrm{Gr}$. No dimensionless
number above depends on an invented constant: groups whose inputs are still datasheet-conditional
are listed with that dependency in Table~\ref{tab:baseline}.
